\documentclass[acmtog, authorversion]{acmart}
\usepackage{graphicx}
%\usepackage[numbers]{natbib}
%\usepackage{float}
%\usepackage{hyperref}
%\usepackage{comment}
%\usepackage{ragged2e}
\setcopyright{none}
\citestyle{acmauthoryear}


\begin{document}

\title{Járművek trajektóriájinak előrejelzése machine learning modellekkel}

\author{Péter Bence Mérnökinformatika BSc 6. félév}%\\Mérnökinformatika BSc 6. félév\\Konzulensek:\\Dr. Horváth András egyetemi docens\\Agg Áron PhD hallgató}
\authornotemark[1]
\affiliation{
    \institution{Széchenyi István Egyetem}
    \city{Győr}
    \country{Hungary}}

\author{Dr. Horváth András}
\affiliation{
    \institution{Széchenyi István Egyetem}
    \city{Győr}
    \country{Hungary}}

\author{Agg Áron PhD hallgató}
\affiliation{
    \institution{Széchenyi István Egyetem}
    \city{Győr}
    \country{Hungary}}

\begin{teaserfigure}
\includegraphics[width=1\columnwidth]{sze_givk_logo.png}
\Description{egyetem_kar}
\end{teaserfigure}

\begin{abstract}
    Az ITS (intelligent transportation system) egyre nagyobb teret 
    hódít napjainkban és rengeteg különböző területen alkalmazzák 
    ezeket a rendszereket. A közlekedési csomópontok elemzése egy 
    frekventált terület az ITS alkalmazásában.
    Célunk, gépi látás és gépi tanulás felhasználásával, közlekedési
    csomópontok elemzésének automatizálása és felgyorsítása. A kutatásban
    lefektetett alapgondolatokat, kifejlesztett keretrendszert és 
    a felmerülő probélmák megoldásait, a gyakorlatban balesetek megelőzésére,
    renitens viselke-dések kiszűrésére és forgalomirányító renszerek támogatásá-ra
    lehet használni.  
    A kutatásban egy trajektória osztályozó módszert ismertetünk, amely objektumdetektálás 
    és objektumkövetés segítségé-vel elemezi a közlekedési 
    csomópontokban elhaladó járművek mozgá-sát. A mozgásuk alapján 
    klaszterezi a trajektóriákat, majd gépi tanulás 
    segítségével predikciót ad az újonnan belépő járművek kilépési 
    pontjára. A módszerhez 6 különböző közlekedési csomópont-ban 
    készített saját videó adatbázisunkat használtuk fel.
    A tesztelt klaszterezési mód-szerek közül (OPTICS, BIRCH, KMeans, DBSCAN)
    az OPTICS algoritmus bizonyult legjobbnak trajektórák klaszterezésére.
    Összehasonlí-tottunk több különböző klasszifikációs módszert 
    a legpontosabb predikció eléréséhez, amelyek: KNN, SVM, GP, DT, 
    GNB, MLP, SGD. A tanul-mányban bemutatott eljárások közül az 
    SVM adta a legpontosabb 90\%-os eredményt.
    %Ezt a pontosságot valós idejű futás közben 5 fps mellett érte el.
    %Ebből azt a következtetést lehet levonni, hogy jobb sebesség elérése
    %érdekében, vagy a feature vectorok dimenzióját kell csökkenteni, vagy
    %érdemes neurális hálót alkalmazni a klasszifikációhoz.
    %A forráskód megtalálható ebben a github repositoriban 
    %\url{https://www.github.com/Pecneb/computer_vision_research}
\end{abstract}

\maketitle

\tableofcontents

\section{Bevezetés}
A városok növekedése egyre nagyobb forgalomhoz vezet, ami a balesetek, forgalmi dugók számát növeli és a levegő minősége is romlik.
Az ITS (intelligent transportation system) fejlesztése a városokban erre megoldást jelenthet. Ez magába foglalja az információs és
kommunikációs technológiák, mint pélául szenzorok, kamerák, kommunikációs hálózatok és adat elemzés fejlesztését. 5G hálózatokon
keresztül, ezek a technológiák összeköthetők a közlekedési eszközökkel. Ehhez okos forgalomirányítási rendszerek kifejlesztésére
van szükség, amik információval tudnak szolgáni a járművekbe szerelt informatikai rendszereknek.
A legértékesebb információt a közlekedésben részvevő járművek jelen és jövőbeli pozíciója jelenti. Pontos és gyors trajektória 
előrejelző rendszerek kifejlesztése egy nagy kihívás és egyre növekszik irántuk a kereslet. E kutatási terület kiforratlanságából
eredően, kevés létező keretrendszer és adathalmaz található, így a tanító adathalmaz gyűjtése, adatok kinyerésének formátuma, tárolása
és mérőszámok kifejlesztése (amivel a tesztelni kívánt modellek pontosságát tudjuk mérni) is a kutatáshoz tartoznak.  
Ebben a kutatásban erre a problémára törekszünk egy módszertant és keretrendszert kifejleszteni, emellett klaszterezési és klasszifikációs
algoritmusokat tesztelni. A tanító adatok előállításához, objektumok detektálására a YOLOv7 \cite{wang2022yolov7}
konvolúciós neurális hálót használtuk, ez a konvolúciósl neurális háló architektúra nem csak nagy pontosságot hanem sebességet is nyújt nekünk. 
Emellett képkockáról képkockára követni is kell tudni a detektált objektumokat. Erre is sok megoldás található manapság, erre a feladatra
a DeepSORT \cite{Wojke2018deep} nevezetű algoritmust használtuk, ez kálmán filtert és konvolúciós neurális hálót használ az objektumok követésére.
A tanító adatok 6 különböző helyszín forgalmát tartalmazzák. Minden helyszín más tulajdonságokkal bír, ezért nem lehet generalizálni
a tanítási folyamatot, nem lehet egy univerzális modellt betanítani ami minden közlekedési helyszínre alkalmazható egyaránt. A klaszterezés
során megpróbáljuk minél pontosabban meghatározni a be és kimeneti pontok által leírt klasztereket, amelyek majd alapul szolgálnak a
klasszifikáció tanítása során. Több fajta klaszterezési algoritmust megvizsgáltunk a kutatás során, KMeans, OPTICS \cite{10.1145/304181.304187}, BIRCH \cite{10.1145/233269.233324} és DBSCAN \cite{10.5555/3001460.3001507}\cite{10.1145/3068335}. 
A klasszifikációhoz bináris klasszifikációs modelleket kombinálunk, így több klasszos klasszifikációs modelt
kapunk. Minden bináris modelnél, egy klassz az összes többivel szemben van betanítva. A modellek pontosságának kiértékelésére 3 mérőszámot
alkalmaztunk, amik az \emph{Accuracy Score}, \emph{Balanced Accuracy Score} \cite{10.1109/ICPR.2010.764} és \emph{Top-k Accuracy Score}.
Mindegyik mérőszám kiszámolásához \emph{K-Fold Cross-Validation} \cite{Anguita2012TheI} metódust alkalmaztunk, ahol \begin{math}K=5\end{math}.

\section{Kapcsolódó kutatások}
Sok ITS-el kapcsolatos kutatásban tárgyalják a forgalom folyás (traffic flow) előrejelzését. \cite{PAUL2017177} össze-hasonlítja az eddig
kutatott és használt modellek, mint például Kalman Filtering, k-nearest neighbor (k-NN), mesterséges neurális hálók, stb., pontosságát és
sebességét, ezen modellek tovább-kutatását, mivel egyre növekednek a különböző szenzorok által  begyűjtött traffic flow adatok, így ez a terület belépett
a \emph{Big Data} korszakába. \cite{10.1371/journal.pone.0253868} is a traffic flow előrejelzését és generálását tárgyalja, Floating Car Data (FCD) 
adathalmazokon betanított, Hosszú-Rövid-Távú memóriájú és Generatív versengő hálókkal.

\bibliographystyle{ACM-Reference-Format}
\bibliography{biblio}

\begin{CCSXML}
    <ccs2012>
        <concept>
            <concept_id>10010147.10010257.10010293</concept_id>
            <concept_desc>Computing methodologies~Machine learning approaches</concept_desc>
            <concept_significance>500</concept_significance>
            </concept>
    </ccs2012>
\end{CCSXML}
\ccsdesc[500]{Computing methodologies~Machine learning approaches}
\end{document}